\documentclass[10pt,a4paper]{amsart}
\usepackage{amsmath,amssymb,amsthm,graphicx}
\usepackage{verbatim}
\newenvironment{reference}{\comment}{\endcomment}
%\newenvironment{reference}{}{}
\newenvironment{slogan}{\comment}{\endcomment}
\newenvironment{history}{\comment}{\endcomment}

% For commutative diagrams we use Xy-pic
\usepackage[all]{xy}

% We use 2cell for 2-commutative diagrams.
\xyoption{2cell}
% Theorem environments.
%
\theoremstyle{plain}
\newtheorem{theorem}[subsection]{Theorem}
\newtheorem{proposition}[subsection]{Proposition}
\newtheorem{lemma}[subsection]{Lemma}

\theoremstyle{definition}
\newtheorem{definition}[subsection]{Definition}
\newtheorem{example}[subsection]{Example}
\newtheorem{exercise}[subsection]{Exercise}
\newtheorem{situation}[subsection]{Situation}

\theoremstyle{remark}
\newtheorem{remark}[subsection]{Remark}
\newtheorem{remarks}[subsection]{Remarks}

\numberwithin{equation}{subsection}

% Macros
%
\def\lim{\mathop{\mathrm{lim}}\nolimits}
\def\colim{\mathop{\mathrm{colim}}\nolimits}
\def\Spec{\mathop{\mathrm{Spec}}}
\def\Hom{\mathop{\mathrm{Hom}}\nolimits}
\def\Ext{\mathop{\mathrm{Ext}}\nolimits}
\def\SheafHom{\mathop{\mathcal{H}\!\mathit{om}}\nolimits}
\def\SheafExt{\mathop{\mathcal{E}\!\mathit{xt}}\nolimits}
\def\Sch{\mathit{Sch}}
\def\Mor{\mathop{\mathrm{Mor}}\nolimits}
\def\Ob{\mathop{\mathrm{Ob}}\nolimits}
\def\Sh{\mathop{\mathit{Sh}}\nolimits}
\def\NL{\mathop{N\!L}\nolimits}
\def\CH{\mathop{\mathrm{CH}}\nolimits}
\def\proetale{{pro\text{-}\acute{e}tale}}
\def\etale{{\acute{e}tale}}
\def\QCoh{\mathit{QCoh}}
\def\Ker{\mathop{\mathrm{Ker}}}
\def\Im{\mathop{\mathrm{Im}}}
\def\Coker{\mathop{\mathrm{Coker}}}
\def\Coim{\mathop{\mathrm{Coim}}}

% Boxtimes
%
\DeclareMathSymbol{\boxtimes}{\mathbin}{AMSa}{"02}

%
% Macros for moduli stacks/spaces
%
\def\QCohstack{\mathcal{QC}\!\mathit{oh}}
\def\Cohstack{\mathcal{C}\!\mathit{oh}}
\def\Spacesstack{\mathcal{S}\!\mathit{paces}}
\def\Quotfunctor{\mathrm{Quot}}
\def\Hilbfunctor{\mathrm{Hilb}}
\def\Curvesstack{\mathcal{C}\!\mathit{urves}}
\def\Polarizedstack{\mathcal{P}\!\mathit{olarized}}
\def\Complexesstack{\mathcal{C}\!\mathit{omplexes}}
% \Pic is the operator that assigns to X its picard group, usage \Pic(X)
% \Picardstack_{X/B} denotes the Picard stack of X over B
% \Picardfunctor_{X/B} denotes the Picard functor of X over B
\def\Pic{\mathop{\mathrm{Pic}}\nolimits}
\def\Picardstack{\mathcal{P}\!\mathit{ic}}
\def\Picardfunctor{\mathrm{Pic}}
\def\Deformationcategory{\mathcal{D}\!\mathit{ef}}
\usepackage{cleveref}
\crefname{originallemma}{original lemma}{original lemmas}
\Crefname{originallemma}{Original lemma}{Original lemmas}
\makeatletter
\newlabel{flat-lemma-bourbaki-finite-type-general-base-at-point}{{05IH}{source}{}{sourcepointer-1}{}}
\newlabel{flat-lemma-bourbaki-finite-type-general-base-at-point@cref}{{[originallemma][0][]05IH}{[1][0][]source}{}{}{}}
\makeatother
\crefname{originaldefinition}{original definition}{original definitions}
\Crefname{originaldefinition}{Original definition}{Original definitions}
\makeatletter
\newlabel{spaces-divisors-definition-relative-weak-assassin}{{0CV0}{source}{}{sourcepointer-2}{}}
\newlabel{spaces-divisors-definition-relative-weak-assassin@cref}{{[originaldefinition][0][]0CV0}{[1][0][]source}{}{}{}}
\makeatother
\crefname{originaldefinition}{original definition}{original definitions}
\Crefname{originaldefinition}{Original definition}{Original definitions}
\makeatletter
\newlabel{spaces-divisors-definition-weakly-associated}{{0CTX}{source}{}{sourcepointer-3}{}}
\newlabel{spaces-divisors-definition-weakly-associated@cref}{{[originaldefinition][0][]0CTX}{[1][0][]source}{}{}{}}
\makeatother
\crefname{originaldefinition}{original definition}{original definitions}
\Crefname{originaldefinition}{Original definition}{Original definitions}
\makeatletter
\newlabel{spaces-morphisms-definition-flat-module}{{05VV}{source}{}{sourcepointer-4}{}}
\newlabel{spaces-morphisms-definition-flat-module@cref}{{[originaldefinition][0][]05VV}{[1][0][]source}{}{}{}}
\makeatother
\title{Pointwise weakly associated criterion for a tensor product over a locally finite type morphism flat at the point}
\author{The Stacks Project Authors}
\date{stacks-project:spaces-flat.tex:local-snapshot-20261012; GNU FDL 1.2 or later}
\begin{document}
\maketitle
\noindent\textit{Editorial note.} All delivered mathematical text is the author's own, copied from the pinned Stacks Project chapter file spaces-flat.tex; no statement or proof step was rewritten, reordered or removed.\par\medskip
\noindent\textit{Editorial note.} The counted proof is the author's complete argument: a surjective etale descent diagram reducing the pointwise weakly associated criterion for algebraic spaces to the scheme case, with the three point conditions translated definition by definition\par\medskip
\noindent\textit{Editorial note.} The label lemma-bourbaki-finite-type-general-base-at-point names the counted statement, and the proof environment immediately following it in the source is delivered verbatim as its proof.\par\medskip
\noindent\textit{Editorial note.} The counted statement is carried by the author's own lemma environment in the chapter source.\par\medskip
\noindent\textit{Allowed standard prerequisites (editorial).}\begin{itemize}
\item The definition of weakly associated points of a quasi-coherent module on an algebraic space, verified after surjective etale descent to schemes
\item The relative assassin Ass\_\{X/Y\}(F) defined through the fibres over the base
\item Flatness of a module at a point of an algebraic space as an etale local condition
\item The scheme case: the Bourbaki criterion for weakly associated points of a tensor product over a locally finite type morphism of schemes flat at the point
\end{itemize}\medskip
\ifdefined\hypertarget\hypertarget{sourcepointer-1}{}\fi
\noindent\textit{Source reference.} Original lemma 05IH: sources/stacks-chow/tags.txt, source lines 7131--7131. Prerequisite pointer invoked at proof line(s) 729; it names a standard fact of a sibling chapter that the counted argument uses, and every step of the counted proof built on it is retained verbatim in the delivered body.\par\medskip
\ifdefined\hypertarget\hypertarget{sourcepointer-2}{}\fi
\noindent\textit{Source reference.} Original definition 0CV0: sources/stacks-chow/tags.txt, source lines 16120--16120. Prerequisite pointer invoked at proof line(s) 723; it names a standard fact of a sibling chapter that the counted argument uses, and every step of the counted proof built on it is retained verbatim in the delivered body.\par\medskip
\ifdefined\hypertarget\hypertarget{sourcepointer-3}{}\fi
\noindent\textit{Source reference.} Original definition 0CTX: sources/stacks-chow/tags.txt, source lines 16082--16082. Prerequisite pointer invoked at proof line(s) 717; it names a standard fact of a sibling chapter that the counted argument uses, and every step of the counted proof built on it is retained verbatim in the delivered body.\par\medskip
\ifdefined\hypertarget\hypertarget{sourcepointer-4}{}\fi
\noindent\textit{Source reference.} Original definition 05VV: sources/stacks-chow/tags.txt, source lines 7564--7564. Prerequisite pointer invoked at proof line(s) 726; it names a standard fact of a sibling chapter that the counted argument uses, and every step of the counted proof built on it is retained verbatim in the delivered body.\par\medskip
\begin{lemma}
\label{lemma-bourbaki-finite-type-general-base-at-point}
Let $S$ be a scheme.
Let $f : X \to Y$ be a morphism of algebraic spaces over $S$ which is
locally of finite type. Let $x \in |X|$ with image $y \in |Y|$.
Let $\mathcal{F}$ be a finite type quasi-coherent sheaf on $X$.
Let $\mathcal{G}$ be a quasi-coherent sheaf on $Y$.
If $\mathcal{F}$ is flat at $x$ over $Y$, then
$$
x \in \text{WeakAss}_X(\mathcal{F} \otimes_{\mathcal{O}_X} f^*\mathcal{G})
\Leftrightarrow
y \in \text{WeakAss}_Y(\mathcal{G})
\text{ and }
x \in \text{Ass}_{X/Y}(\mathcal{F}).
$$
\end{lemma}

\begin{proof}
Choose a commutative diagram
$$
\xymatrix{
U \ar[d] \ar[r]_g & V \ar[d] \\
X \ar[r]^f & Y
}
$$
where $U$ and $V$ are schemes and the vertical arrows are surjective \'etale.
Choose $u \in U$ mapping to $x$. Let $\mathcal{E} = \mathcal{F}|_U$
and $\mathcal{H} = \mathcal{G}|_V$.
Let $v \in V$ be the image of $u$. Then
$x \in \text{WeakAss}_X(\mathcal{F} \otimes_{\mathcal{O}_X} f^*\mathcal{G})$
if and only if
$u \in \text{WeakAss}_X(\mathcal{E} \otimes_{\mathcal{O}_X} g^*\mathcal{H})$
by Divisors on Spaces, Definition
\ref{spaces-divisors-definition-weakly-associated}.
Similarly, $y \in \text{WeakAss}_Y(\mathcal{G})$ if and only if
$v \in \text{WeakAss}_V(\mathcal{H})$.
Finally, we have $x \in \text{Ass}_{X/Y}(\mathcal{F})$ if and only if
$u \in \text{Ass}_{U_v}(\mathcal{E}|_{U_v})$ by
Divisors on Spaces, Definition
\ref{spaces-divisors-definition-relative-weak-assassin}.
Observe that flatness of $\mathcal{F}$ at $x$ is
equivalent to flatness of $\mathcal{E}$ at $u$, see
Morphisms of Spaces, Definition \ref{spaces-morphisms-definition-flat-module}.
The equivalence for $g : U \to V$, $\mathcal{E}$, $\mathcal{H}$, $u$, and $v$
is More on Flatness, Lemma
\ref{flat-lemma-bourbaki-finite-type-general-base-at-point}.
\end{proof}

\end{document}
